\chapter{Fazit und Ausblick}
% Die Schlussfolgerungen beenden den Artikel und erhöhen den Abstraktionsgrad der Beschreibung wieder, genau so, wie die Einleitung diesen gesenkt hat. Meist sind Schlussfolgerungen nicht sehr sorgfältig geschrieben, weil sie unter Termindruck entstanden sind und als nicht wichtig erachtet werden. Dieser Praxis sollte nicht gefolgt werden, da viele Begutachtende nach "Zusammenfassung (Abstract)" und "Einleitung (Introduction)" zunächst die "Schlussfolgerungen (Conclusion)" lesen und dann beurteilen, ob sie den ganzen Artikel lesen. Idealerweise bestehe die Schlussfolgerungen aus folgenden drei Punkten, für die jeweils ein Abschnitt gesetzt werden sollte:
Text
% Zusammenfassung: Was wurde in dieser Arbeit gemacht? Was waren die Schlüsselergebnisse? Diesmal jedoch zusammengefasst nochmals für Lesende, die die vorherigen Seiten mit allen Details bereits gelesen haben.
% Adressaten der Verbesserung: Wem nützen die Verbesserungen/Beiträge der Arbeit? Inwieweit wird der Stand der Wissenschaft und Technik durch die Arbeit verbessert?
% Aufbauende/Zukünftige Arbeiten: Welche nächsten Schritte sind geplant (erst kurzfristige, dann längerfristige)? Welche möglichen Lösungsansätze für noch bestehende Probleme sind denkbar? Wie könnten Folgearbeiten aussehen?

% Kritische Würdigung: Vergleich mit verwandten Arbeiten, Diskussion offener Punkte
% und insbesondere auch Ihre persönliche Bewertung der Arbeit. Was ist gut gelungen, was nicht?


% \begin{table}[ht!]
%     \centering
%     \begin{tabular}{|p{3cm}|p{1cm}|p{5cm}|}
%     \hline
%     \textbf{Symbol} & 
%     \textbf{Unit} &
%     \textbf{Parameter}
%     \\
%     \hline
%     \var $P_{r}$ &
%     $[W]$ &
%     received power
%     \\
%     \hline
%     \var $P_{t}$ &
%     {\ttfamily $[W]$} &
%     transmitted power
%     \\
%     \hline
%     \var $R_{t}$ &
%     {\ttfamily $[m]$} &
%     Tx seperation
%     \\
%     \hline
%     \var $R_{r}$ &
%     {\ttfamily $[m]$} &
%     Rx seperation
%     \\
%     \hline
%     \var \lambda &
%     {\ttfamily $[m]$} &
%     wavelength
%     \\
%     \hline
%     \var \sigma &
%     {\ttfamily $[m^2]$} &
%     radar cross section
%     \\
%     \hline
%     \end{tabular}
%     \caption{parameters of bistatic radar system for point target}
%     \label{table:bistatic radar parameters}
%     \end{table}
    
%     \vspace{3cm}

% %Wenn so günstig warum machen wir das nicht so? 
